\documentclass[11pt]{article}

\usepackage[a4paper,margin=1in]{geometry}
\usepackage[T1]{fontenc}
\usepackage[utf8]{inputenc}
\usepackage{lmodern}
\usepackage{amsmath,amssymb}
\usepackage{physics}
\usepackage{hyperref}
\usepackage{enumitem}
\usepackage{setspace}
\usepackage{mdframed}

\hypersetup{colorlinks=true,linkcolor=blue,citecolor=blue,urlcolor=blue}
\setstretch{1.1}

\newmdenv[linewidth=1.1pt,innertopmargin=5pt,innerbottommargin=5pt,%
  innerleftmargin=8pt,innerrightmargin=8pt]{resultbox}

\title{SSV Conclusions~A\\[2pt]
\large Finite update rate and presentist emergent gravity:\\
what the programme leaves to mainstream physics}
\author{Stig Norland}
\date{Closing note, 2026-06-16}

\begin{document}
\maketitle

\begin{abstract}
The Saturated Superfluid Vacuum (SSV) programme is closed: no surviving route
derives the quantities it set out to derive. This note records what of its
\emph{ontology} maps cleanly onto standing physics, so the conceptual content is
not lost with the framework. The central picture --- local time as the rate of
phase update of an order parameter, a finite update capacity that loading
consumes, and gravity as the response of a holographic record to that loading ---
is a coherent position, but it is not new and it makes no prediction beyond the
programmes it restates. It is the entropic/thermodynamic emergent-gravity
programme (Jacobson, Verlinde, Padmanabhan) on a superfluid vacuum (Volovik,
Zloshchastiev), with time read as a Josephson phase rate and a presentist reading
of the emergent metric. Its honest home is foundations / philosophy of physics:
interpretive coherence, not a gravity theory. Following the repository's
standing rule, the magnitude failures are stated at full strength and not
softened.
\end{abstract}

\section{Purpose and the honest frame}

This is a conclusions document, not a research paper. It exists because the SSV
ontology --- specifically \emph{time delay from limited update capacity} and
\emph{gravitational bending as a holographic-screen update} --- is the part of
the programme that still feels like it might say something, and the owner asked
whether it can be adapted to mainstream physics. The answer below is: the
ontology is real, recognisable, and already named in the literature; what it
adds on top is interpretation rather than prediction; and the reasons it cannot
be more are the same reasons the programme closed. Per the repository's negative
results rule, those reasons are stated plainly.

\section{The ontology in one paragraph}

Physical time is identified with the local rate of change of the vacuum order
parameter $\Psi=\sqrt{\rho}\,e^{i\theta}$. In a condensate the phase obeys a
Josephson relation,
\begin{equation}
  \hbar\,\dot\theta = -\mu ,
\end{equation}
so the local ``tick'' is set by the local chemical potential. The medium has a
\emph{finite} capacity to update its state; loading it (placing energy/defects in
it) consumes update capacity and slows the local tick. Where the slowing is
uneven it has a gradient, and any motion through the gradient is bent toward the
region of greatest delay. That bending is gravity (Paper~IV). In the holographic
reading (Paper~VII-b, the ``Option~D'' fork) nothing bends space in the bulk:
$\Psi$ lives \emph{on} space, the loading is written to a self-defined acoustic
horizon as an area-law record, and the read-back of that record steers
trajectories. Matter writes the ledger; the ledger dents capacity; the dent is
gravity.

\section{Where each component already lives}

\begin{itemize}[leftmargin=1.4em]
\item \textbf{Time as phase rate $\to$ Josephson + analogue $g_{00}$.} Promoting
  $\hbar\dot\theta=-\mu$ to an emergent lapse is standard analogue gravity: the
  acoustic metric's $g_{00}$ is fixed by the local condensate state
  (Unruh~1981; Visser~1998; Volovik). Weak-field, this reproduces
  $g_{00}=1+2\Phi/c^2$ and the Einstein-1911 (time-delay) half of light
  deflection. This is exposition of a known result.
\item \textbf{Finite update capacity $\to$ Margolus--Levitin / Lloyd.} The one
  genuinely distinct gloss --- a finite \emph{rate} of state change --- has a
  rigorous mainstream anchor: the maximum rate of orthogonal-state evolution is
  bounded by energy, $\nu_\perp \le 2E/\pi\hbar$ (Margolus--Levitin~1998;
  Lloyd, \emph{Ultimate physical limits to computation}, 2000). SSV's ``update
  capacity'' is a physical-substrate gloss on this bound. It does not, however,
  turn the bound into a new gravitational prediction.
\item \textbf{Time as change, not backdrop $\to$ relational time.} The claim that
  time is the medium changing, not a dimension it moves through, is the
  relational tradition: Barbour's timeless mechanics and the Connes--Rovelli
  thermal time hypothesis (1994). SSV is a concrete realisation, not a new
  argument.
\item \textbf{Screen update as gravity $\to$ Jacobson + Verlinde + Padmanabhan.}
  ``Gravity is the holographic screen updating'' is, almost verbatim, entropic /
  thermodynamic gravity: Jacobson~1995 (the Einstein equation as an equation of
  state, $\delta Q = T\,dS$), Verlinde~2011 (gravity as an entropic force from
  holographic screens), and Padmanabhan's gravity-as-thermodynamics. The
  Option~D arc runs through exactly this machinery.
\end{itemize}

\section{What the ontology adds, and what it does not}

What it adds is \emph{interpretive coherence}: a single medium in which the same
phase field carries the local clock (rate $\dot\theta$), the static near-field
(gradient $\nabla\theta$), and the long-range encoding (the gapless propagating
mode). That unification is genuinely attractive as a picture.

What it does not add is a falsifiable prediction. The construction reproduces
the \emph{form} of the weak-field metric, the area law, and the Unruh/Hawking
temperature, and concedes every \emph{magnitude}: Newton's constant $G$, the
cosmological scale $\Lambda/H_0$, and the MOND-like acceleration $a_0$ all enter
as inputs rather than outputs. ``Form yes, magnitude no'' is not an incidental
gap; it is the wall the programme hit. Conclusions~B records the sharpest
instance of it.

\section{An internal tension the salvage must own}

The two halves of the ontology are, on inspection, two different ontologies:

\begin{itemize}[leftmargin=1.4em]
\item \emph{Update capacity (Paper~IV)} is \textbf{bulk-realist}: the medium's
  local state physically slows time and the gradient bends motion.
\item \emph{Holographic screen (Option~D)} is \textbf{anti-bulk-realist}:
  ``nothing bends space in the bulk; space was never bulk-dynamical; geometry is
  localised into the encoding.''
\end{itemize}

\noindent These reconcile only through a \textbf{duality} claim --- bulk and
screen are two complete descriptions of one state, with no independent dials.
That duality is the most rhetorically powerful element of the picture and the one
with the least empirical content: by the programme's own audit it ``carries no
falsifier of its own.'' Any honest adaptation should present the duality as a
heuristic frame (the Wheeler/Deser ``laws as closure conditions'' lineage), not
as a mechanism.

\section{Presentism: the one place the ontology pulls its weight}

SSV is natively presentist. There is local becoming --- the update tick, a
privileged substrate ``now'' --- but no completed, global, infinite time and no
block universe. This is legitimate precisely \emph{because} SSV's spacetime is
emergent rather than fundamental: an effective metric with a domain of validity
is under no obligation to be geodesically complete, so there is no real
future infinity $\mathcal{I}^+$ for a global construction to require. This is not
idle: it dissolves the obstacle that breaks global de~Sitter holography
(the non-unitary Euclidean CFT at $\mathcal{I}^+$ has no substrate referent) and
routes naturally to \emph{observer / static-patch} holography, where well-defined
observables and a finite entropy require including the observer's worldline ---
the ``now'' --- as in the de~Sitter observer-algebra programme
(Chandrasekaran--Longo--Penington--Witten~2022). The caveat, stated plainly: the
present cosmological horizon still has a Gibbons--Hawking entropy whose
microstates SSV does not supply, so presentism removes an obstacle without
building the dual.

\section{Verdict}

\begin{resultbox}
\noindent The ontology --- finite-update-rate time dilation plus
holographic-screen gravity --- is a coherent foundations position that restates
the entropic/emergent-gravity programme on a presentist superfluid vacuum. It is
adaptable to mainstream physics as \textbf{philosophy of physics}: a careful
essay on time as finite update rate and on presentist emergent gravity, engaging
Jacobson, Verlinde, Padmanabhan, Barbour, Connes--Rovelli, and the de~Sitter
observer-algebra work. It is \emph{not} adaptable as a new theory of gravity: it
predicts nothing those programmes do not, it inherits their open empirical
objections (e.g.\ the qBounce and interferometric constraints on entropic
gravity), and in concrete form it fails on magnitude.
\end{resultbox}

\paragraph{Do not reclaim under a new banner.} The closed no-gos remain closed:
deriving $G$, a spin-2 graviton or tensor gravitational-wave polarisations, a
fresh resolution of Loschmidt's paradox, and $a_0$ as a selected prediction.
These are recorded as falsified or conceded and should not be re-imported into a
foundations write-up as live claims.

\begin{thebibliography}{9}
\bibitem{josephson} B.\,D.\ Josephson, \emph{Possible new effects in
  superconductive tunnelling}, Phys.\ Lett.\ \textbf{1}, 251 (1962).
\bibitem{unruh} W.\,G.\ Unruh, \emph{Experimental black-hole evaporation?},
  Phys.\ Rev.\ Lett.\ \textbf{46}, 1351 (1981).
\bibitem{visser} M.\ Visser, \emph{Acoustic black holes: horizons, ergospheres
  and Hawking radiation}, Class.\ Quantum Grav.\ \textbf{15}, 1767 (1998);
  gr-qc/9712010.
\bibitem{volovik} G.\,E.\ Volovik, \emph{The Universe in a Helium Droplet}
  (Oxford, 2003).
\bibitem{zlosh} K.\,G.\ Zloshchastiev, \emph{Logarithmic nonlinearity in
  theories of quantum gravity: origin of time and observational consequences},
  Grav.\ Cosmol.\ \textbf{16}, 288 (2010).
\bibitem{margolus} N.\ Margolus and L.\,B.\ Levitin, \emph{The maximum speed of
  dynamical evolution}, Physica D \textbf{120}, 188 (1998).
\bibitem{lloyd} S.\ Lloyd, \emph{Ultimate physical limits to computation},
  Nature \textbf{406}, 1047 (2000).
\bibitem{barbour} J.\ Barbour, \emph{The End of Time} (Oxford, 1999).
\bibitem{connesrovelli} A.\ Connes and C.\ Rovelli, \emph{Von Neumann algebra
  automorphisms and time--thermodynamics relation in general covariant quantum
  theories}, Class.\ Quantum Grav.\ \textbf{11}, 2899 (1994).
\bibitem{jacobson} T.\ Jacobson, \emph{Thermodynamics of spacetime: the Einstein
  equation of state}, Phys.\ Rev.\ Lett.\ \textbf{75}, 1260 (1995).
\bibitem{verlinde} E.\ Verlinde, \emph{On the origin of gravity and the laws of
  Newton}, JHEP \textbf{04}, 029 (2011).
\bibitem{padmanabhan} T.\ Padmanabhan, \emph{Thermodynamical aspects of gravity:
  new insights}, Rep.\ Prog.\ Phys.\ \textbf{73}, 046901 (2010).
\bibitem{clpw} V.\ Chandrasekaran, R.\ Longo, G.\ Penington, E.\ Witten,
  \emph{An algebra of observables for de~Sitter space}, JHEP \textbf{02}, 082
  (2023); arXiv:2206.10780.
\end{thebibliography}

\end{document}
